% GENERATED by scripts/make_proxy_tables.py -- do not edit by hand.
\begin{table}[t]
\centering
\caption{Tahsis uygunluğu (öncelik-ağırlıklı zamanında tamamlama), 500 tohum, 5 robot / 25 görev. \textbf{Sol:} mükemmel-navigasyon düzlemi, tahsis kalitesini yalıtır. \textbf{Sağ:} stokastik navigasyon vekili; burada Nav2 hedef arızası da bir göreve mal olabilir. AHE-MRTA her iki düzlemde ve her senaryoda öndedir; temel yöntemlere karşı eşli Wilcoxon karşılaştırmalarının biri dışında hepsi anlamlıdır (arıza senaryosunda AHE--BiG beraberedir). \textbf{Kalın} sütun en iyisidir.}
\label{tab:fitness}
\small
\setlength{\tabcolsep}{4pt}
\begin{tabular}{l ccc ccc}
\toprule
& \multicolumn{3}{c}{\textbf{Mükemmel nav (tahsis kalitesi)}} & \multicolumn{3}{c}{\textbf{Stokastik nav (dayanıklılık)}} \\
\cmidrule(lr){2-4}\cmidrule(lr){5-7}
\textbf{Yöntem} & RF & MS & DP & RF & MS & DP \\
\midrule
\textbf{AHE-MRTA*} & \textbf{0.994} & \textbf{0.727} & \textbf{0.840} & \textbf{0.384} & \textbf{0.263} & \textbf{0.349} \\
BiG-MRTA & 0.980 & 0.668 & 0.733 & 0.380 & 0.229 & 0.284 \\
RoSTAM-EA & 0.935 & 0.643 & 0.679 & 0.333 & 0.225 & 0.227 \\
Cons-DBTA & 0.978 & 0.642 & 0.729 & 0.364 & 0.217 & 0.258 \\
\bottomrule
\end{tabular}
\end{table}
